% ============================================================================
\chapter{Constrained Learning of Volatility Surfaces}\label{ch:learning}
% ============================================================================

The previous chapter ended with a family for which the no-arbitrage conditions are
\emph{theorems}: choose SSVI parameters inside the Gatheral--Jacquier region and the surface is
provably clean, on the whole domain, forever. The price of that guarantee is a five-parameter
functional form, and Section~\ref{sec:parametric} will show what it costs in fit.

This chapter turns to model classes that pay no such price and enjoy no such guarantee. A neural
network can represent essentially any smooth surface, so the rigidity disappears; but no theorem
attaches to it, and the arbitrage conditions must instead be imposed \emph{during training}, by
penalising violations wherever the training procedure happens to look. That last clause is
innocuous-sounding and is the subject of the entire chapter.

Section~\ref{sec:softtheory} develops the theory of that approximation and is the theoretical
backbone of the thesis: what minimising a penalty to zero actually proves, how far the statement
extends beyond the points where it was checked, why the butterfly and calendar constraints
demand incompatible sampling grids, and when a measured violation is large enough to matter
economically. Sections~\ref{sec:nnsec} and~\ref{sec:opsec} then supply the
approximation-theoretic background for the two learned families studied empirically, the per-day
neural smoother and the neural operator, the second of which changes the object being learned
from a surface to the map that produces surfaces. Section~\ref{sec:problem} closes by stating
the reconstruction problem the rest of the thesis solves and scores, including the one criterion
by which our framing departs from the literature.

% ============================================================================
\section{Soft constraints, sampling sets, and materiality}\label{sec:softtheory}
% ============================================================================

Parametric certificates are theorems; neural smoothers have none. They cannot impose
$g\ge0$ and $\partial_\tau w\ge0$ pointwise, and instead penalise violations on a finite
\emph{collocation set} $\hat X$ \cite{ackerer2020,hoshisashi2023}:
\begin{equation}\label{eq:loss}
\mathcal L=\underbrace{\tfrac1p\textstyle\sum_i\omega_i\bigl(w_\theta(k_i,\tau_i)-w_i\bigr)^2}_{\text{fit}}
+\lambda_{\mathrm b}\,\overline{\relu(-g_\theta)^2}\Big|_{\hat X_{\mathrm b}}
+\lambda_{\mathrm c}\,\overline{\relu(-\partial_\tau w_\theta)^2}\Big|_{\hat X_{\mathrm c}}
+\lambda_{\mathrm w}\,\overline{(\partial_{kk}w_\theta)^2}\Big|_{\hat X_{\mathrm{wing}}},
\end{equation}
where overlines denote averages over the indicated collocation grids and $\omega_i$ are fit
weights. This section builds the small amount of theory the thesis needs about
\eqref{eq:loss}: what minimising a penalty to zero \emph{proves}, how far that statement
extends between nodes, why the two arbitrage constraints have incompatible grid requirements,
and when a measured violation is economically meaningful.

\begin{remark}[A penalty is a statement about $\hat X$]\label{rem:sampling}
Minimising \eqref{eq:loss} to $\lambda$-term zero proves $g_\theta\ge0$ \emph{on}
$\hat X_{\mathrm b}$ and $\partial_\tau w_\theta\ge0$ \emph{on} $\hat X_{\mathrm c}$. Whether that extends between the nodes depends on the interaction between the
network's inductive bias and the geometry of the constraint.
Sections~\ref{sec:deep}--\ref{sec:operator} show empirically that it can fail, exactly, and at
scale. The methodological consequence adopted throughout is an \emph{independent audit grid}:
finer than, and disjoint from, every collocation grid, with all penalties reported twice (at
the nodes and on the audit) and a \emph{blind-spot flag} raised when the two disagree
materially.
\end{remark}

\subsection{The node-gap bound}

The following elementary proposition formalises the mechanism: between nodes, the only control
a node-wise constraint gives is through a regularity bound on the model class, and the
violation budget scales \emph{linearly} with the node gap.

\begin{theorem}[Node-gap violation bound]\label{prop:nodegap}
Fix $k$ and let $f(\tau):=\partial_\tau w_\theta(k,\tau)$ on an interval containing adjacent
collocation nodes $\tau_1<\tau_2$ with gap $h=\tau_2-\tau_1$. Suppose $f(\tau_1)\ge0$,
$f(\tau_2)\ge0$ (the penalty is satisfied at the nodes) and $f$ is Lipschitz with
$|f'(\tau)|=|\partial_{\tau\tau}w_\theta(k,\tau)|\le L$ on $[\tau_1,\tau_2]$. Then
\[
f(\tau)\;\ge\;-\,\frac{L\,h}{2}\qquad\text{for all }\tau\in[\tau_1,\tau_2],
\]
and the bound is sharp: for every $L,h$ there is an $f$ with $|f'|\le L$, $f(\tau_{1,2})=0$
and $\min f=-Lh/2$. Consequently, the deepest calendar violation a node-satisfying surface can
hide between nodes is proportional to the node gap, with constant the curvature budget of the
model class.
\end{theorem}

\begin{proof}
For $\tau\in[\tau_1,\tau_2]$, let $\tau^\ast\in\{\tau_1,\tau_2\}$ be the nearer node, so
$|\tau-\tau^\ast|\le h/2$. By the Lipschitz bound,
$f(\tau)\ge f(\tau^\ast)-L|\tau-\tau^\ast|\ge-Lh/2$. Sharpness: take
$f(\tau)=L\,|\tau-\tfrac{\tau_1+\tau_2}{2}|-\tfrac{Lh}{2}$ (or a $C^1$ mollification thereof),
which vanishes at both nodes, has $|f'|=L$, and attains $-Lh/2$ at the midpoint.
\end{proof}

\begin{remark}[Reading the bound]\label{rem:nodegapread}
Two consequences drive Section~\ref{sec:deep}. First, \emph{without} an a-priori bound $L$ on
$\partial_{\tau\tau}w_\theta$, node constraints control nothing between nodes, and a trained
network under fit pressure has no useful such bound: the fit term of \eqref{eq:loss} is then
free to buy accuracy with a dip of any depth, at zero penalty cost. Second, \emph{with}
whatever $L$ the trained network does exhibit, the violation budget is proportional to the node
gap $h$. Shrinking the calendar grid's maximal gap from $0.41$y, the exponential grid's
long-end hole, to $0.017$y, the repaired uniform grid, shrinks the worst admissible hidden
violation by a factor of about $24$ at equal curvature. This is exactly the repair adopted in
Section~\ref{sec:deep}, and the reason the repaired node and audit penalties agree to grid
noise. Figure~\ref{fig:nodegap} illustrates the extremal $f$, and Figure~\ref{fig:gridhole}
shows the two grids on the thesis's actual maturity axis.
\end{remark}

\thfig{figures/fig_nodegap.pdf}{0.78}{The node-gap bound
(Theorem~\ref{prop:nodegap}): a surface can satisfy the calendar penalty exactly at two
adjacent collocation nodes (dots, $\partial_\tau w=0$) while $\partial_\tau w$ dips to
$-Lh/2$ between them (red area). The worst-case violation is linear in the gap $h$ and in the
curvature budget $L$.}{fig:nodegap}

\thfig{figures/fig_grid_hole.pdf}{1.0}{Constraint-specific collocation on the thesis's
maturity axis $[0.06,1.4]$. The literature's exponential grid (top, ten nodes,
$0.060,0.085,\dots,0.987,1.400$) is optimal for the butterfly constraint but leaves a
$0.413$-year hole, $30\%$ of the axis (shaded), exactly where the calendar stressor of
Section~\ref{sec:deep} forces a crossing. The calendar constraint therefore gets its own
uniform grid (bottom, 80 nodes, maximal gap $0.017$y).}{fig:gridhole}

\subsection{Why one grid cannot serve both constraints}

The two constraints of Definition~\ref{def:durrleman} and Definition~\ref{def:calendar} have
\emph{opposite} grid requirements:


\begin{center}
\small
\begin{tabularx}{\textwidth}{@{}lXX@{}}
\toprule
\textbf{Constraint} & \textbf{Where it bites} & \textbf{Required grid in $\tau$} \\
\midrule
Butterfly $g\ge0$
& Short end, where small $w$ amplifies the term $\tfrac{(w')^2}{4w}$ in \eqref{eq:g}
& Dense near the short end, for example an exponential grid \\

Calendar $\partial_\tau w\ge0$
& Anywhere along the maturity axis, since it involves a $\tau$-derivative
& Uniformly bounded node gaps, by Theorem~\ref{prop:nodegap} \\
\bottomrule
\end{tabularx}
\end{center}

The butterfly constraint is most restrictive when total variance is small because the slope term in \eqref{eq:g}, amplified by $1/w$, dominates in that region. This also explains why NB02 violations concentrate between $7$ and $14$ days. An exponential maturity grid therefore places more nodes near the short end.

The calendar constraint is different. It is a first-derivative condition that may fail anywhere along the maturity axis. Theorem~\ref{prop:nodegap} shows that its enforcement quality depends on the largest node gap. An exponential grid creates this largest gap at the long end. Using the same grid for both constraints therefore weakens the calendar penalty at long maturities. This is the caveat identified by Chataigner et al.~\cite{chataigner2020} and the mechanism behind the collocation blind spot discussed in Section~\ref{sec:deep}.


\subsection{Regularity of the model class: activations}
The butterfly penalty needs $w''_\theta$, so the corrector's activation function must itself be
twice differentiable, and with a curvature the training can actually shape. Recall the two
canonical choices:
\begin{equation}\label{eq:activations}
\relu(x)=\max(x,0),
\qquad
\tanh(x)=\frac{e^{x}-e^{-x}}{e^{x}+e^{-x}} .
\end{equation}
Their second derivatives could not be more different. The rectifier is piecewise affine, so
$\relu''(x)=0$ for every $x\neq0$ and is undefined at the single kink; a network built from it
is continuous piecewise-affine \cite{hornik1991,baydin2018}, and therefore has $\partial_{kk}\mathrm{MLP}_\theta=0$ Lebesgue-almost
everywhere. The hyperbolic tangent is $C^\infty$ with $\tanh''(x)=-2\tanh(x)\bigl(1-\tanh^2(x)\bigr)$,
a smooth non-degenerate curvature at every point.

The consequence for \eqref{eq:g} is immediate. With a $\relu$ corrector, the $w''$ channel of
$g_\theta$ carries no information at almost every collocation node: there the corrector
contributes zero curvature, $w''_\theta$ reduces to the second derivative of the smooth SSVI
prior alone, and the butterfly penalty cannot reshape it. The penalty's gradient reaches the
corrector only through the level and slope terms of $g$, never through curvature, which is
precisely the quantity the constraint exists to control.

The requirement is therefore $C^2$ regularity with non-degenerate curvature, not $\tanh$
specifically, and it is worth being explicit that modern practice supplies several admissible
alternatives. GELU and SELU, which have largely displaced $\tanh$ as defaults in contemporary
architectures, are both smooth and both satisfy the requirement; the neural operators of
Section~\ref{sec:opsec} use GELU internally for that reason. We retain $\tanh$ for the per-day
corrector of \eqref{eq:arch} on a narrower ground: its bounded range keeps the multiplicative
correction numerically controlled near initialisation, where training starts at the prior and we
want the corrector to depart from it gradually. The content of this subsection is the exclusion
rather than the choice. Piecewise-affine activations make the butterfly penalty structurally
blind to the very derivative it constrains, and no amount of tuning elsewhere in the training
recovers that.

% \begin{remark}[Exactness of the nested derivatives]\label{rem:c2}
% The derivatives entering the penalties are computed by automatic differentiation and validated
% against closed forms: in $k$ against the SVI derivatives \eqref{eq:svideriv} (maximum error
% $5.6\times10^{-17}$ and $1.1\times10^{-16}$ for $w'$ and $w''$), and in $\tau$ against central
% differences on the SSVI prior ($1.4\times10^{-11}$). The $\tau$ path is validated separately
% because the calendar penalty depends on nothing else. \notebookfig{03}{AD vs closed form and
% finite differences}
% \end{remark}

\thfig{figures/fig_relu_tanh.pdf}{1.0}{Curvature as seen by the butterfly penalty,
toy illustration: a piecewise-affine (ReLU-type) fit of a convex smile
has $w''=0$ almost everywhere with distributional spikes at the knots (red), so the penalty's
curvature channel is inert; a $\tanh$-type fit exposes a smooth, informative
$w''$ (blue).}{fig:relutanh}

\subsection{Materiality}
An audit that flags every sign change of $\partial_\tau w$ or $g$, however small, is useless: a
soft-constraint method leaves tiny numerical residuals everywhere, and calling them
``violations'' would drown the real ones. We therefore need a threshold below which a violation
is indistinguishable from quote noise and above which it is economically meaningful. That
threshold is the materiality tolerance.

\begin{definition}[Materiality threshold]\label{def:materiality}
A violation is \emph{material} if it exceeds $\mathrm{tol}_{\mathrm{mat}}=10^{-3}$ in surface
units: $w$-units for $g$-type quantities (through their effect on $w$), and
$\partial_\tau w$-units for calendar.
\end{definition}

The value $10^{-3}$ is not arbitrary; it is calibrated to the noise in the quotes themselves.
Multiplicative implied-volatility noise of relative size $\varepsilon$ propagates to total
variance by the delta method, $\delta w/w = 2\,\delta\sigma/\sigma \approx 2\varepsilon$. With
$\varepsilon\approx0.5$--$0.75\%$ (half the $1$--$1.5\%$ multiplicative band observed on $w$) and
$w\approx0.03$ at the long end, the noise band on $w$ is about $4\times10^{-4}$. A calendar
violation must therefore sustain $|\partial_\tau w|\gtrsim10^{-3}$ over half a year to move $w$
beyond this band. Below that, a ``violation'' is a numerical residual of the soft-constraint
approach, not a monetisable opportunity; Hoshisashi et al.\ make the same observation for their
DCNN \cite{hoshisashi2023}.

Two consequences hold throughout the thesis. Every violation rate is reported twice, \emph{strict}
($<-10^{-8}$, any sign change at all) and \emph{material} ($<-\mathrm{tol}_{\mathrm{mat}}$,
economically real), so the reader can always separate the two. And blind-spot detection compares
the node grid and the audit grid \emph{on the same material scale}: an earlier,
magnitude-inconsistent detector false-alarmed on residuals four orders of magnitude below
materiality (Section~\ref{sec:deep}), which the material-scale comparison eliminates.

\subsection{The audit-grid methodology}
Remark~\ref{rem:sampling} and Theorem~\ref{prop:nodegap} together say that a soft penalty
certifies its constraint only on the finite set where it is sampled, and that the violation it
can hide between samples grows with the gap. A penalty reporting zero is therefore not evidence
of an arbitrage-free surface; it is evidence only about the sampling set. Everything downstream
in this thesis rests on refusing to take a penalty's own report at face value, and instead
auditing each model against a standard it was not trained on. Concretely, four rules are applied
throughout.

\begin{enumerate}[label=(A\arabic*)]
\item \textbf{Two grids, never one.} Every penalty is evaluated twice: once on its own
collocation grid, which is what the training loss sees and minimises, and once on an
\emph{independent audit grid} that is finer than, and disjoint from, every collocation grid.
The training grid measures how well the constraint was \emph{imposed}; the audit grid measures
whether it actually \emph{holds}. Only their disagreement is informative, which is why the two
must not share nodes.

\item \textbf{Compare on the material scale.} Node and audit statistics are compared at the
materiality tolerance of Definition~\ref{def:materiality}, not in raw units, so that numerical
residuals are not mistaken for arbitrage. A \emph{blind-spot flag} is raised precisely when the
nodes are materially clean while the audit is materially violating, i.e.\ when the training loss
would report success on a surface that an honest audit rejects.

\item \textbf{For operators, audit unseen days too.} A per-day smoother is audited between its
nodes; an operator must additionally be audited on \emph{days} disjoint from training, because
its penalty is a statement about training-day surfaces only, with nothing guaranteeing that the
constraint transports to a new day (Section~\ref{sec:operator}). The blind spot thus acquires a
second axis, space and time, for operators.

\item \textbf{Gate every stress and invariance claim.} Before an arbitrage comparison on a
stressed or subsampled input is read, a fit-validity gate certifies that the model actually
fitted that input. Without the gate, a model that simply failed to fit could be mistaken for one
that fitted cleanly, and a collapsed operator's trivially perfect invariance could be reported
as a genuine property (Sections~\ref{sec:deep} and~\ref{sec:operator}).
\end{enumerate}

% ============================================================================
\section{Neural networks as surface models}\label{sec:nnsec}
% ============================================================================

The parametric models of Sections~\ref{sec:svisec}--\ref{sec:ssvisec} impose a functional form
and let the data choose five parameters. Their strength is that the form carries theorems; their
weakness is that the form is also a ceiling, since no choice of $(a,b,\rho,m,\sigma)$ can produce
a shape SVI cannot express. This section introduces the opposite trade: a model class flexible
enough to represent essentially any smooth surface, whose price is that it carries no theorem at
all and must be constrained by the penalties of Section~\ref{sec:softtheory}.

\subsection{From parametric forms to learned functions}

A neural network is best understood as a \emph{parameterised family of functions} rich enough to
approximate anything, rather than as a black box. The basic building block, the
\emph{multilayer perceptron} (MLP), alternates affine maps and a fixed nonlinearity. With
activation $\phi_a:\R\to\R$, depth $D$ and widths $n_0,\dots,n_D$, it is the map
$x\mapsto z_D(x)$ defined recursively by
\begin{equation}\label{eq:mlp}
z_0=x,\qquad
z_\ell=\phi_a\bigl(W_\ell z_{\ell-1}+b_\ell\bigr)\quad(\ell<D),\qquad
z_D=W_Dz_{D-1}+b_D,
\end{equation}
with parameters $\theta=(W_\ell,b_\ell)_{\ell\le D}$ and $\phi_a$ applied componentwise. In our
setting the input is a point of the surface, $x=(k,\tau)$, and the output is a scalar. Figure
\ref{fig:mlp} shows the object.

\begin{figure}[htbp]\centering
\begin{tikzpicture}[
  neuron/.style={circle,draw=black!70,minimum size=5.5mm,inner sep=0pt},
  nodein/.style={circle,draw=black!70,fill=black!8,minimum size=5.5mm,inner sep=0pt},
  nodeout/.style={circle,draw=black!70,fill=black!8,minimum size=5.5mm,inner sep=0pt},
  arr/.style={-{Stealth[length=1.4mm]},black!30,line width=0.25pt},
  lbl/.style={font=\scriptsize}
]
\foreach \i/\nm in {1/{$k$},2/{$\tau$}}
  \node[nodein] (I\i) at (0,{1.1-1.0*\i}) {\scriptsize\nm};
\foreach \i in {1,...,4}
  \node[neuron] (H\i) at (2.1,{2.0-0.85*\i}) {};
\foreach \i in {1,...,4}
  \node[neuron] (G\i) at (4.2,{2.0-0.85*\i}) {};
\node[nodeout] (O) at (6.3,-0.4) {};
\foreach \i in {1,2}\foreach \j in {1,...,4} \draw[arr] (I\i)--(H\j);
\foreach \i in {1,...,4}\foreach \j in {1,...,4} \draw[arr] (H\i)--(G\j);
\foreach \i in {1,...,4} \draw[arr] (G\i)--(O);
\node[lbl,align=center] at (0,1.0) {input\\$(k,\tau)$};
\node[lbl,align=center] at (2.1,1.75) {hidden\\$\phi_a(W_1x+b_1)$};
\node[lbl,align=center] at (4.2,1.75) {hidden\\$\phi_a(W_2z_1+b_2)$};
\node[lbl,align=center] at (6.3,1.0) {output\\$\mathrm{MLP}_\theta(k,\tau)$};
\end{tikzpicture}
\caption{A multilayer perceptron \eqref{eq:mlp}: it maps \emph{one point} of the surface to
\emph{one number}. Each arrow is a learned weight; each circle applies the activation $\phi_a$.
To obtain a whole surface, the network is evaluated point by point, and the parameters $\theta$
are fitted to one day's quotes.}\label{fig:mlp}
\end{figure}

The classical justification for this class is that it is dense in the space of continuous
functions, and, crucially for us, dense in derivatives too.

\begin{proposition}[{Universal approximation; Hornik \cite{hornik1991}}]\label{thm:hornik}
If $\phi_a$ is continuous, bounded and non-constant, then for every compact $K\subset\R^d$ the
set of one-hidden-layer MLPs is dense in $C(K)$ for the uniform norm. Moreover, if
$\phi_a\in C^m$, density holds in the $C^m$ norm on compacts: the network approximates the
target \emph{and} its derivatives up to order $m$ simultaneously.
\end{proposition}

The $C^m$ statement is the relevant one here, and not a technicality. The quantities penalised
in \eqref{eq:loss} are $w''_\theta$ and $\partial_\tau w_\theta$, so approximation power must
reach the derivatives, not only the level; and by the curvature argument of
Section~\ref{sec:softtheory} the activation must be at least $C^2$ for the butterfly channel to
be trainable at all.

\subsection{Exact derivatives by automatic differentiation}

A constrained smoother is only as good as its derivatives, and this subsection explains why a
learned smoother can have exact ones at all. The idea is not native to finance: it is the
mechanism underlying physics-informed neural networks \cite{raissi2019,lagaris1998}, where a
differential equation is imposed by penalising its residual at sampled points, which is only
meaningful because the residual can be evaluated exactly. Automatic differentiation (AD)
computes derivatives \emph{exactly}, to floating-point precision, rather than by finite
differences:
the program defining $w_\theta$ is a composition of elementary operations whose local
derivatives are known in closed form, and the chain rule is applied mechanically along the
computational graph \cite{baydin2018}.

Two consequences matter. First, the surface derivatives entering \eqref{eq:loss}, namely
$w'_\theta$ and $w''_\theta$ in $k$ and $\partial_\tau w_\theta$ in $\tau$, are exact functionals
of $\theta$; the penalties are therefore differentiable in $\theta$ and their gradients are
themselves exact, a nested application of AD. Second, exactness is \emph{testable}: the AD
derivatives can be checked against the closed-form SVI derivatives \eqref{eq:svideriv} in $k$ and
against central differences in $\tau$, which is how a silent wiring error in a derivative path
would be caught. The neural operators of Section~\ref{sec:opsec}, following the reference design
\cite{wiedemann2025}, instead resolve their penalties by finite differences on a grid; the gap
between what that grid sees and what an independent audit sees is itself one of this thesis's
measurements.

\subsection{The multiplicative prior--corrector architecture}

Used alone, an MLP would discard everything the parametric theory provides. The architecture of
\cite{ackerer2020}, adopted in Section~\ref{sec:deep}, instead keeps the certified surface and
learns only the deviation from it:
\begin{equation}\label{eq:arch}
w_\theta(k,\tau)=\underbrace{w_{\mathrm{SSVI}}(k,\tau)}_{\text{certified prior}}\times
\underbrace{\exp\bigl(\mathrm{MLP}_\theta(\tilde k,\tilde\tau)\bigr)}_{\text{positive corrector}},
\end{equation}
where the SSVI prior is fitted to the day's quotes and certified arbitrage-free by
Theorem~\ref{lem:range}, and the corrector is a $\tanh$ MLP on rescaled inputs, initialised so
that the exponential is $\approx1$ and training therefore \emph{starts at the prior}.

Three properties of \eqref{eq:arch} are used repeatedly. \emph{Positivity} of $w_\theta$ holds by
construction, since a product of a positive prior and an exponential cannot vanish or change
sign. \emph{Additivity in log-space}, $\log w_\theta=\log w_{\mathrm{SSVI}}+\mathrm{MLP}_\theta$,
means the corrector models the \emph{relative} deviation from the prior, which also explains a
failure mode observed in Section~\ref{sec:deep}: a wrong prior is worse than a flat one, because
the corrector must first spend capacity undoing it. \emph{Inheritance}: wherever the corrector is
flat, all derivatives of $w_\theta$, hence $g$ and $\partial_\tau w$, are the prior's, and the
prior carries a theorem.

% ============================================================================
\section{Neural operators}\label{sec:opsec}
% ============================================================================

\subsection{What is being learned: functions versus operators}

The models above learn a \emph{function}: given the numbers $(k,\tau)$, return a number. Their
parameters are fitted to one day's quotes, and a new day requires a new optimisation from
scratch. Neural operators change the object being learned. Instead of one surface, they learn
the \emph{map that produces surfaces}.

It is worth being precise about the word \emph{map}, since the whole distinction rests on it. A
map is simply a rule assigning to each element of one set exactly one element of another. What
varies is what those sets contain. A \emph{function}, in the ordinary sense, maps numbers to
numbers, for instance $(k,\tau)\mapsto w$. An \emph{operator} maps \emph{functions to functions}:
its input is not a number but an entire object, and so is its output. Differentiation is the
canonical example, sending a function $f$ to the function $f'$; so is integration. In our
problem, the operator of interest is
\begin{equation}\label{eq:opmap}
\mathcal G:\;\underbrace{a}_{\text{a day's quote cloud}}\;\longmapsto\;
\underbrace{w(\cdot,\cdot)}_{\text{the whole smoothed surface}},
\qquad
a=\bigl\{(k_i,\tau_i,\sigma_{\mathrm{IV},i})\bigr\}_{i\le p},
\end{equation}
which takes an entire configuration of quotes and returns an entire surface, evaluable at any
$(k,\tau)$ we care to query. Formally $a$ is a discrete measure, that is, a discretisation of an
element of a function space, and $\mathcal G$ acts between function spaces. Figure~\ref{fig:funcvsop}
contrasts the two levels.

\begin{figure}[htbp]\centering
\begin{tikzpicture}[
  bx/.style={draw=black!70,rounded corners=1pt,minimum width=20mm,minimum height=8mm,
              font=\scriptsize,align=center},
  net/.style={draw=black!70,fill=black!6,rounded corners=1pt,minimum width=17mm,
              minimum height=9mm,font=\scriptsize,align=center},
  ar/.style={-{Stealth[length=1.6mm]},black!75},
  capt/.style={font=\scriptsize\itshape,align=center}
]
\node[bx] (fin) at (0,0) {one point\\$(k,\tau)$};
\node[net] (fnet) at (3.1,0) {MLP $_\theta$};
\node[bx] (fout) at (6.2,0) {one number\\$w(k,\tau)$};
\draw[ar] (fin)--(fnet); \draw[ar] (fnet)--(fout);
\node[capt] at (-2.15,0) {function\\(a network)};
\node[bx] (oin) at (0,-2.0) {a whole cloud\\$a=\{(k_i,\tau_i,\sigma_i)\}$};
\node[net] (onet) at (3.1,-2.0) {$\mathcal G_\vartheta$};
\node[bx] (oout) at (6.2,-2.0) {a whole surface\\$w(\cdot,\cdot)$};
\draw[ar] (oin)--(onet); \draw[ar] (onet)--(oout);
\node[capt] at (-2.15,-2.0) {operator\\(a neural operator)};
\draw[black!25,dashed] (-3.3,-1.0)--(7.5,-1.0);
\end{tikzpicture}
\caption{The two levels of learning. Top: a network maps \emph{numbers to numbers}; its
parameters encode one surface, and a new day requires refitting. Bottom: an operator maps a
\emph{whole quote configuration to a whole surface}; its parameters encode the smoothing
\emph{procedure}, so a new day is a single forward pass.}\label{fig:funcvsop}
\end{figure}

The practical payoff is amortisation. Training is paid once, across the entire history, after
which a new day costs one forward pass, milliseconds rather than the minute or so needed to
refit a per-day smoother. Two structural properties distinguish a genuine operator from a plain
network fed with a flattened input. It must accept inputs of \emph{any size and layout}, since
the number of quotes and their locations change every day, which requires invariance under
permutation of the index $i$ and tolerance of variable $p$. And it aspires to
\emph{discretization invariance}: refining the sampling of the input should not change the
answer.

\begin{definition}[Discretization invariance, after \cite{kovachki2023}]\label{def:discinv}
A parametric map $\mathcal G_\vartheta$ on point-cloud inputs is discretization-invariant at $u$
if, for any sequence of discretisations $a^{(n)}$ of $u$ refining to $u$, the outputs
$\mathcal G_\vartheta(a^{(n)})$ converge to a common limit $\mathcal G_\vartheta(u)$, uniformly
on compacts. Empirically one tests a finite version: the output surface must be stable under
subsampling of the input quotes. Section~\ref{sec:operator} gates this test on fit validity,
because a collapsed operator that outputs a constant is trivially invariant.
\end{definition}

\subsection{DeepONet: a learned global basis}

The first architecture answers the question ``how can a network output an entire function?'' by
writing that function in a \emph{basis}. DeepONet \cite{lu2021} represents the operator in
separated form,
\begin{equation}\label{eq:deeponet}
\mathcal G_\vartheta(a)(k,\tau)\;=\;\sum_{j=1}^{p^\ast}
\underbrace{b_j\bigl(E(a)\bigr)}_{\text{coefficients: depend on the quotes}}\;
\underbrace{t_j(k,\tau)}_{\text{basis functions: fixed}}\;+\;b_0 .
\end{equation}
Three pieces do the work. The \emph{encoder} $E$ compresses the irregular quote cloud into a
fixed-length vector and must be permutation-invariant, for which we use the Deep Sets form
\begin{equation}\label{eq:deepsets}
E(a)=\rho_E\Bigl(\sum_{i\le p}\phi_E(k_i,\tau_i,\sigma_i)\Bigr),
\end{equation}
the general shape of a continuous permutation-invariant map \cite{zaheer2017}: the inner sum
discards the ordering of the quotes, and the outer network restores expressiveness. The
\emph{branch} network $b:\R^{d_E}\to\R^{p^\ast}$ turns that summary into $p^\ast$ coefficients.
The \emph{trunk} network $t:\R^2\to\R^{p^\ast}$ produces $p^\ast$ smooth functions of the query
point, and it does \emph{not} see the quotes at all. The day's information therefore enters only
through the coefficients, while the shapes available are fixed once training is over.

\begin{figure}[htbp]\centering
\begin{tikzpicture}[
  blk2/.style={draw=black!70,fill=black!4,rounded corners=1pt,minimum width=18mm,
               minimum height=8mm,font=\scriptsize,align=center},
  blk/.style={draw=black!70,rounded corners=1pt,minimum width=18mm,minimum height=8mm,
              font=\scriptsize,align=center},
  net/.style={draw=black!70,fill=black!6,rounded corners=1pt,minimum width=16mm,
              minimum height=8mm,font=\scriptsize,align=center},
  ar/.style={-{Stealth[length=1.6mm]},black!75},
  ann/.style={font=\scriptsize\itshape,align=center,text=black!60}
]
\node[blk] (cloud) at (0,0.9) {quote cloud\\$a$};
\node[net] (enc)   at (2.6,0.9) {encoder $E$\\(perm.-inv.)};
\node[net] (br)    at (5.2,0.9) {branch $b$};
\node[blk] (coef)  at (7.9,0.9) {coefficients\\$b_1,\dots,b_{p^\ast}$};
\node[blk] (query) at (0,-1.2) {query point\\$(k,\tau)$};
\node[net] (tr)    at (2.6,-1.2) {trunk $t$};
\node[blk] (basis) at (5.2,-1.2) {basis\\$t_1,\dots,t_{p^\ast}$};
\node[blk2] (sum) at (7.9,-1.2) {$\sum_j b_j t_j+b_0$\\$=w(k,\tau)$};
\draw[ar] (cloud)--(enc); \draw[ar] (enc)--(br); \draw[ar] (br)--(coef);
\draw[ar] (query)--(tr);  \draw[ar] (tr)--(basis); \draw[ar] (basis)--(sum);
\draw[ar] (coef)--(sum);
\node[ann] at (3.9,-2.35) {the trunk never sees the quotes: the basis is global and smooth};
\end{tikzpicture}
\caption{DeepONet \eqref{eq:deeponet}. The quotes determine only the \emph{coefficients} (top
path); the \emph{shapes} come from a trunk basis that is the same for every day (bottom path).
The output is their weighted sum.}\label{fig:deeponet}
\end{figure}

The construction is backed by an operator version of universal approximation.

\begin{proposition}[{Chen--Chen \cite{chenchen1995}; Lu et al.\ \cite{lu2021}}]\label{thm:deeponet}
Let $\mathcal G$ be a continuous operator from a compact subset of a Banach space of functions to
$C(K)$, with $K\subset\R^d$ compact. Then for every $\varepsilon>0$ there exist networks $b,t$,
and an encoder acting on point evaluations of the input, such that the separated form
\eqref{eq:deeponet} approximates $\mathcal G$ uniformly to within $\varepsilon$.
\end{proposition}

The structural consequence exploited in Sections~\ref{sec:operator} and~\ref{ch:downstream} is
visible directly in \eqref{eq:deeponet}: for \emph{every} input, the output lies in the span of
the same $p^\ast$ trunk functions, a global, input-independent, smooth basis. Whatever the quotes
do, the surface inherits the trunk's regularity, and fine-scale curvature is simply not
expressible. This is why the DeepONet turns out to be largely immune to butterfly violations,
its residual weakness being a thin calendar blind spot, at the price of visibly coarser fits.

\subsection{Graph neural operators: a learned local kernel}

The second architecture answers the same question differently, by building the output
\emph{locally} from nearby data rather than from a global basis. The neural-operator programme \cite{li2020gno,li2021fno,kovachki2023} models $\mathcal G$ as a
composition of \emph{kernel integral layers},
\begin{equation}\label{eq:kernelint}
v_{\ell+1}(x)\;=\;\phi_a\Bigl(W_\ell\,v_\ell(x)\;+\;\int\kappa_\ell(x,y)\,v_\ell(y)\,
\nu(\dd y)\Bigr),
\end{equation}
with learned kernels $\kappa_\ell$ and $\nu$ the input measure. The integral is what makes this
an operator rather than a network: the value at $x$ depends on the input \emph{everywhere},
weighted by a learned kernel, instead of on a fixed-length vector of features.

On an irregular quote cloud the integral cannot be computed exactly, so it is discretised by
\emph{message passing} over a graph. For each query point $x$ one keeps its $m$ nearest quotes
$N_m(x)$ and forms
\begin{equation}\label{eq:gno}
\int\kappa_\ell(x,y)\,v_\ell(y)\,\nu(\dd y)\;\approx\;
\frac{1}{m}\sum_{y\in N_m(x)}\kappa_\ell\bigl(x,y,a(y)\bigr)\,v_\ell(y),
\end{equation}
a Monte-Carlo quadrature of \eqref{eq:kernelint} localised to the neighbourhood. Operator Deep
Smoothing \cite{wiedemann2025} instantiates exactly this design for implied volatility and is
the reference for the graph-neural-operator (GNO) variant of Section~\ref{sec:operator}.

\begin{figure}[htbp]\centering
\begin{tikzpicture}[
  quo/.style={circle,fill=black!65,inner sep=0pt,minimum size=1.6mm},
  qpt/.style={circle,draw=black!80,fill=white,line width=0.5pt,inner sep=0pt,minimum size=2.6mm},
  msg/.style={-{Stealth[length=1.3mm]},black!55,line width=0.4pt},
  ann/.style={font=\scriptsize,align=center},
  blk/.style={draw=black!70,rounded corners=1pt,minimum width=16mm,minimum height=8mm,
              font=\scriptsize,align=center}
]
\foreach \a/\b in {0.30/1.35, 0.95/1.75, 1.60/1.20, 0.55/0.55, 1.35/0.35,
                   2.35/1.55, 2.75/0.75, 0.15/0.30, 2.05/0.15, 1.95/2.05,
                   3.25/1.30, 2.60/2.20}
  \node[quo] at (\a,\b) {};
\node[qpt] (X) at (1.30,1.10) {};
\draw[black!35,dashed] (1.30,1.10) circle (0.72);
\foreach \a/\b in {0.95/1.75, 1.60/1.20, 0.55/0.55, 1.35/0.35}
  \draw[msg] (\a,\b)--(X);
\node[ann] at (1.30,-0.35) {query $x$ and its $m$ nearest quotes $N_m(x)$};
\node[blk] (out) at (5.6,1.10) {$v_{\ell+1}(x)$\\{\scriptsize local kernel sum \eqref{eq:gno}}};
\draw[-{Stealth[length=1.6mm]},black!75] (2.55,1.10)--(out);
\node[ann,text=black!60] at (5.6,-0.35) {each query is assembled\\from its own neighbourhood};
\end{tikzpicture}
\caption{Graph neural operator \eqref{eq:gno}. The surface value at a query point $x$ is
assembled by message passing from the $m$ nearest quotes, a discretisation of the kernel
integral \eqref{eq:kernelint}. Neighbourhoods change as $x$ moves, which buys local adaptivity
and, as Remark~\ref{rem:gnobias} explains, costs global control of curvature.}\label{fig:gno}
\end{figure}

\begin{remark}[Inductive bias and the audit]\label{rem:gnobias}
The locality of \eqref{eq:gno} is simultaneously the GNO's strength and its liability. The
strength is adaptivity: the output at $x$ is assembled from the nearest quotes, so the operator
follows local smile features closely and leads the accuracy table of Section~\ref{sec:operator}.
The liability is that nothing global constrains the assembled surface's second derivative.
Neighbourhoods change discretely as $x$ moves and each carries its own kernel response, so the
output exhibits fine-scale curvature at the scale of the quote spacing. A butterfly penalty
resolved by finite differences on a coarse grid, here the $22\times9$ grid of the reference
design with maximal $\tau$-gap $0.176$y, is blind to curvature at finer scales in exactly the
sense of Theorem~\ref{prop:nodegap}; an independent $61\times31$ audit is not. The empirical
counterpart, that the accuracy leader is also the most arbitrage-dirty model in the study, is
Section~\ref{sec:operator}'s central exhibit.
\end{remark}

\subsection{Why this matters here}

The two architectures embody opposite answers to the same question, and the contrast organises
the empirical chapters. DeepONet fixes a smooth global basis and lets the data choose
coefficients, so regularity is guaranteed by construction and fine detail is unreachable. The
graph operator assembles each query from its own neighbourhood, so fine detail is reachable and
regularity is guaranteed nowhere. Table~\ref{tab:opcompare} summarises the comparison, including
the per-day smoother of Section~\ref{sec:nnsec} for reference.

\begin{table}[htbp]\centering\small
\caption{The three model classes compared. The last column is what
Chapters~\ref{ch:methods}--\ref{ch:downstream} measure.}\label{tab:opcompare}
\begin{tabular}{lccc}
\toprule
 & per-day MLP smoother & DeepONet & graph neural operator \\
\midrule
learns a & function (one surface) & operator & operator \\
new day costs & a full refit & one forward pass & one forward pass \\
output shapes & free & span of a global basis & locally assembled \\
regularity & from prior and penalty & by construction & none imposed \\
derivatives for penalties & exact (AD) & exact (AD) & finite differences \\
expected failure mode & between nodes & coarse fits & fine-scale curvature \\
\bottomrule
\end{tabular}
\end{table}

Both operator families are trained with soft penalties, and both therefore inherit the question
of Section~\ref{sec:softtheory}: on which set is the constraint actually enforced? For a per-day
smoother the answer concerned the space between collocation nodes. For an operator a second gap
opens, because the penalty is minimised on \emph{training days} and nothing transports that
property to a day the model has never seen. Section~\ref{sec:operator} measures both gaps.

% ============================================================================
\section{The reconstruction problem}\label{sec:problem}
% ============================================================================
The pieces are now in place to state precisely the problem the rest of the thesis solves and
scores, and to make explicit the one criterion by which its framing departs from the literature.

Given a single day's quotes $\{(k_i,\tau_i,\sigma_{\mathrm{IV},i})\}_{i=1}^{p}$, sparse,
irregular and noisy, the task is to construct a surface
$\widehat w:\R\times(0,\infty)\to(0,\infty)$ that

\begin{enumerate}[label=(\roman*)]
\item \textbf{fits} the observed quotes;
\item \textbf{is arbitrage-free}, i.e.\ satisfies Definitions~\ref{def:durrleman}
and~\ref{def:calendar} and hence, by Lemma~\ref{lem:gjbutterfly} and
Proposition~\ref{prop:calendar}, admits no static arbitrage;
\item \textbf{generalises} to quotes held out from the fit; and, the criterion this thesis adds,
\item \textbf{is honest about its own guarantee}: it states the set on which (ii) is actually
enforced, and the economic size of any residual violation off that set.
\end{enumerate}

Criteria (i)--(iii) are the standard trio, and every method in
Chapters~\ref{ch:methods}--\ref{ch:downstream} is scored on them identically. Criterion (iv) is
the one the thesis argues cannot be omitted. The reason is the whole content of
Section~\ref{sec:softtheory}: for every model class except the certified parametric ones, (ii)
is imposed only on a finite sampling set, and a method can satisfy (i)--(iii) beautifully while
failing (ii) between its samples, or, for an operator, on days it never saw. A surface that
scores well on the classical trio is therefore not yet trustworthy; what makes it trustworthy is
knowing \emph{where} its no-arbitrage guarantee holds and \emph{what it costs} where it does not.
A model that fits, excludes arbitrage on a stated set, generalises, and prices its residual
failure is doing something a model that merely fits, generalises, and reports a zero penalty is
not.

Each criterion is scored by a specific instrument, all shared across families. The unified
protocol of Chapter~\ref{ch:data} scores (i)--(iii). The audit-grid methodology (A1)--(A4) of
Section~\ref{sec:softtheory}, together with the materiality calculus of
Definition~\ref{def:materiality}, scores the guarantee half of (iv). And
Chapter~\ref{ch:downstream} scores its cost half, turning a violation into a repair rate, a
negative density mass, or a degraded signal through the identity
$\sigma_{\mathrm{loc}}^2=\partial_\tau w/g$ of Proposition~\ref{prop:dupire}. With the problem
and its scoring fixed, the benchmark can begin.